\section{Method}
\label{sec:method}

\subsection{Clean Feature Anchoring}
\label{sec:anchor}

Let $\Phi_t$ and $h_t$ be the frozen feature map and classifier of the natural teacher. The student
inherits both, $\Phi_s\leftarrow\Phi_t$ and $h_s=h_t$, after which only $\Phi_s$ is optimized. For a
clean example $x$, CFA solves
\begin{equation}
\begin{aligned}
x_{\mathrm{adv}}
&=\operatorname*{arg\,max}_{x'\in B(x,\epsilon)}
  \lVert\Phi_s(x')-\Phi_t(x)\rVert_2^2,\\
\mathcal L_{\mathrm{CFA}}
&=\mathbb E_x\!\left[\lVert\Phi_s(x_{\mathrm{adv}})-\Phi_t(x)\rVert_2^2\right].
\end{aligned}
\label{eqn:anchor}
\end{equation}
We approximate the inner maximization with randomly initialized PGD against the fixed target
$\Phi_t(x)$ and minimize the same discrepancy over the backbone. This adversarial-training stage
uses no label loss; inference uses $h_t\circ\Phi_s$. \Cref{alg:cfa} gives the complete procedure.

\subsection{Sensitivity-Matched $\epsilon$}
\label{sec:angeps}

A uniform pixel radius need not produce a uniform change in the feature-space objective. We therefore
allocate the batch budget according to each example's clean-point loss sensitivity. For
$\ell_{\mathrm{anc}}(x';x_i)=\lVert\Phi_s(x')-\Phi_t(x_i)\rVert_2^2$, define
\begin{equation}
g_i=\left.\lVert\nabla_{x'}\ell_{\mathrm{anc}}(x';x_i)\rVert_1\right|_{x'=x_i},
\qquad
\max_{\lVert\delta\rVert_\infty\le e}
\langle\nabla_{x'}\ell_{\mathrm{anc}},\delta\rangle=e g_i.
\label{eqn:firstorder}
\end{equation}
The linearized loss change under radius $\epsilon_i$ is therefore $\epsilon_i g_i$. For $g_i>0$,
equalizing this quantity while preserving $\sum_i\epsilon_i=N\epsilon$ gives
\begin{equation}
\epsilon_i=N\epsilon\frac{g_i^{-1}}{\sum_jg_j^{-1}},
\qquad\text{or more generally}\qquad
\epsilon_i\propto g_i^{-p},
\label{eqn:angeps_ideal}
\end{equation}
where $p=0$ recovers a uniform radius. We use $p=1$ and an $\ell_2$ gradient norm as a sensitivity
proxy, measured with one backward pass at clean inputs. The $\ell_1$ version implied by
\Cref{eqn:firstorder} gives $0.29$ points lower AA and retains the clean-accuracy gain
(\Cref{app:box}); exact first-order equalization is therefore a motivation for the implemented rule.

We floor the proxy at $10^{-12}$, clip its inverse-mean multipliers to $[0.5,1.5]$, and normalize
them to unit mean before setting $\epsilon_i$. The floor gives a uniform allocation when all clean
gradients vanish at teacher initialization. Normalization preserves $\sum_i\epsilon_i=N\epsilon$,
although the final multipliers can leave the clipping interval. PGD step sizes are scaled by the same
multipliers. \Cref{app:box} specifies this procedure and an optional exact-box variant.
The matched-budget comparison in \Cref{tab:allocation} tests whether the sensitivity signal improves
on uniform, difficulty-based, and margin-based allocation.
